\chapter{Construct by Construct}\label{ch:constructs}

Every protocol has faithful rules (Chapter~\ref{ch:faithful}), but how they look, and what
they cost, depends on the constructs the protocol is built from. One number tells: the
\emph{lookahead}, how many messages ahead you must look to know where the conversation is.
At lookahead $0$ the last message says it, and every rule is one transition row, like the
rows of Images. Each step of lookahead nests a name one \texttt{someNext}/\texttt{allNext}
level deeper.

Here is each construct's smallest protocol, as LeanFM measures it:

\begin{table}[ht]
\centering
\ToolOutput{generated/ctl-constructs.tex}
\caption{Each construct's smallest protocol: checker nodes (states before merging),
typestates, lookahead, and whether its typestates are deterministic
(Section~\ref{def:det}).}
\label{tab:constructs}
\end{table}

and here is what each one does to the state machine, in one line:

\begin{center}
\begin{tabular}{@{}lp{.62\textwidth}@{}}
\toprule
construct & what it adds \\
\midrule
sequence & one typestate per message \\
choice & one typestate per branch, and a rule naming every branch \\
loop & nothing beyond its body: the next round reuses the same typestates \\
bounded loop & a hidden counter, paid for in lookahead \\
parallel & every combination of where the branches are: sizes multiply \\
group & nothing: brackets for people, invisible to the checks \\
sessions & a log of many sessions splits back into one run per session \\
where & nothing in the control; the data is checked on the run \\
\bottomrule
\end{tabular}
\end{center}

\section{Sequence and choice}\label{prop:answered}

A sequence contributes one typestate per message. A choice contributes one typestate per
first message of a branch, and its row lists every branch. As long as no message appears at
two points that behave differently, the last message names the typestate, and the lookahead
is $0$. A \texttt{cause} adds nothing to the order of messages: it records which request an
answer answers, which matters for logs (Chapter~\ref{ch:where}).

Properties of a protocol are consequences, not extra axioms. Take ``every request is
answered'': after a request $q$, one of its answers $r_1$ or $r_2$ is sure to come.
\begin{lstlisting}[language={},basicstyle=\ttfamily\footnotesize]
(always (implies (on q) (inevitably (or (on r1) (on r2)))))
\end{lstlisting}
If it is true of the protocol, it is true of every system that obeys the protocol's rules,
because the rules say everything the protocol says (Section~\ref{cor:best}). Nobody has to
state it, test it or remember it. Chapter~\ref{ch:mcp} shows LeanFM proposing such
properties and deciding each.

\section{Loops}

A loop plays its body again from the start. The moment after the last message of one round
and the moment after the same message in the next round have the same future, so they are
the same typestate: a loop adds none beyond its body's. The polling protocol of
Chapter~\ref{ch:global} has five checker nodes and four typestates; two nodes merge because
they behave alike.

\section{Bounded loops: a hidden counter}

A bounded loop is unrolled into one copy of its body per try, so it carries a counter the
messages do not show. In Retry, three typestates are reached by Ask, one per try, and three
by Busy. They differ only in how many tries remain, and that shows only in what comes
next: after the third Busy comes GiveUp, not Ask.

\begin{figure}[ht]
\centering
\AccessibleFigure{Typestate graph of Retry: start then Ask then Ok or Busy; after Busy Ask again up to three times; after the third Busy GiveUp. Ok and GiveUp have double borders.}{\Typestates{ctl-retry}}
\caption{Retry's nine typestates. Three are reached by Ask and three by Busy; only the future tells them apart.}
\label{fig:retry-ts}
\end{figure}

So Retry has lookahead $4$, and its names have to describe four steps of future. Written
out as a tree, a name repeats every part it uses at every branch, so it roughly doubles with
each step. Written with shared parts, each part appears once:

\ToolOutput{generated/ctl-retry-sizes.tex}

\noindent Here is the start typestate's name with its shared parts labelled
(\texttt{\#3=} defines a part, \texttt{\#3\#} uses it again):

\ToolOutput{generated/ctl-retry-start.tex}

\noindent and the same name read out by LeanFM, each shared part once:

\ToolOutput{generated/ctl-retry-start-english.tex}

The rule of thumb: a counter that counts to $c$ costs a lookahead of about $c$. Written as
trees the names grow exponentially in $c$; with shared parts they grow by a few lines per
step. Chapter~\ref{ch:deciding} builds the checker on that, and Chapter~\ref{ch:faithful}
counts the same cost in bits: putting the counter in the message makes the lookahead $0$.

\section{Parallel composition}

\texttt{parallel} lets its branches interleave, so a state has to record where
\emph{each} branch is. The last message says only which branch moved last, and nothing
about the others.

\begin{figure}[ht]
\centering
\AccessibleFigure{Typestate graph of Both: two independent request and answer pairs interleaved in every order thirteen typestates.}{\Typestates{ctl-parallel}}
\caption{Two independent requests interleaved: thirteen typestates, lookahead 1. This is the graph on the cover.}
\label{fig:parallel-ts}
\end{figure}

Two request--answer pairs have three typestates each alone, and thirteen interleaved, with
lookahead $1$. In general the sizes multiply: $b$ branches of $t$ typestates each can give
up to $t \times t \times \cdots \times t = t^b$. This is where size really comes from in
practice: GoServer's 1{,}604 typestates come from polling the game list beside the lobby,
and from a seat's four child sessions running side by side.

\section{Brackets: group}

A response is tied to its request by \texttt{cause}, and the view folds the pair into one
arrow. Some runs of messages have no such tie: the frames that carry one HTTP message, the
steps of a handshake. Where to bracket them is a choice, so the protocol writes it:
\begin{lstlisting}[language={},basicstyle=\ttfamily\footnotesize]
protocol Upload (Client, Server, Log) {
  group Request {
    q: Headers from Client to Server;
    Data from Client to Server;
    End from Client to Server;
  }
  Saved from Server to Client cause q;
  parallel {
    group Audit {
      Note from Server to Log;
      Ack from Log to Server;
    }
  } and {
    Done from Server to Client;
  }
}
\end{lstlisting}
A group is for people only. Every check sees its steps in place, so the typestates, rules
and verdicts are those of the same text without the brackets. A bracket can always be
drawn, but folding it into one arrow claims its steps happen together, and that is true
only where nothing else can run between them. \texttt{Request} folds to one arrow from
Client to Server. \texttt{Audit} sits in a \texttt{parallel} branch, where \texttt{Done}
can land between \texttt{Note} and \texttt{Ack}, so it is drawn open and never folds. The
same holds inside \texttt{sessions}, and after a \texttt{join} of fewer than all of its
branches.

\section{Sessions decompose a product}\label{prop:sessions}

A product is large because one model must remember every branch at once. A log does not
have to. Each message carries its session's identity (Chapter~\ref{ch:where}), so a log of
interleaved sessions splits into one run per session, and each run is checked against its
own protocol:

\begin{quote}
A log of independent sessions (no message belongs to two of them) follows their
composition exactly when each session, pulled out of the log by its identity, follows its
own protocol.
\end{quote}

Why: independent sessions share no messages and may interleave in any order, so the
composition allows exactly the logs that shuffle together one allowed run of each.
Pulling a session out by its identity undoes the shuffle. So the cost is the sum of the
sessions' sizes, not their product. Checked one session at a time, every GoServer
protocol has lookahead $0$ except the seat, which keeps its own \texttt{parallel}. A join,
a session that belongs to two others, is not independent, and LeanFM accepts a join only
where the text writes one.

\section{Determinism and prev}\label{def:det}\label{thm:walk}

The typestates of a protocol are \emph{deterministic} when, from any typestate, a given
message leads to only one typestate. Then where the conversation is after a message
depends only on where it was before and which message came, and that is exactly what a
message carrying a \texttt{prev} pointer provides on the wire. Following the chain forward,
each message's typestate is worked out from the one before.

The same idea checks a whole candidate system in one pass. Walk the candidate from its
start, pair each of its states with the protocol's typestate it should be, and check at
each pair that the two have the same label and the same moves:

\begin{lstlisting}[language={},basicstyle=\ttfamily\footnotesize]
def walk(candidate, protocol):             # is the candidate faithful?
    pair = {candidate.start: protocol.start}
    queue = [candidate.start]
    while queue:
        v = queue.pop(0); C = pair[v]
        if label(v) != label(C):
            return False                   # not the same moment
        if {label(w) for w in next(v)} != {label(D) for D in next(C)}:
            return False                   # a move more, or a move less
        for w in next(v):
            D = step(C, label(w))          # only one: the typestates are deterministic
            if pair.get(w, D) != D:
                return False               # one state, two different moments
            if w not in pair:
                pair[w] = D; queue.append(w)
    return True
\end{lstlisting}

Each step of the candidate is looked at once. If the walk succeeds, every reachable state
of the candidate is paired with a typestate that has its label and its moves, step after
step: the candidate has the protocol's moves, which is what faithful means
(Section~\ref{thm:faithful}). If the candidate does have the protocol's moves, no check can
fail, because determinism leaves the walk one typestate to choose at each step, and it is
the right one.
When the typestates are not deterministic, as in GoServer, where two interleaved sessions
can each produce the same message, LeanFM decides faithfulness by refining both models
together instead (Chapter~\ref{ch:deciding}).

\section{Where clauses in the reduction}

A where clause is about values, and the rules are about the order of messages, so a where
clause adds nothing to the rules; Chapter~\ref{ch:where} checks it on the run. Two kinds
still change what the rules are applied to.

\begin{itemize}
\item A clause that \emph{correlates} (\texttt{i.id = g.id}, an echoed request id)
  defines identity. It decides which events form one session, and so which runs the rules
  are checked against (Section~\ref{prop:sessions}).
\item A clause that \emph{guards} a branch (\texttt{o.total >= 100}) decides, on a given
  run, which of the allowed branches may be taken. The row before the choice still names
  both branches as possible; the guard is checked on the values.
\end{itemize}

Either way the protocol's control is plain CTL, and the data is checked where the data is.
